\documentclass[11pt,a4paper]{article}
\usepackage[margin=2.5cm]{geometry}
\usepackage{amsmath,amssymb,graphicx,booktabs,hyperref,caption}
\usepackage{times}
\newcommand{\Wmat}{\mathbf{W}}\newcommand{\Jmat}{\mathbf{J}}
\newcommand{\Xmat}{\mathbf{X}}\newcommand{\Zmat}{\mathbf{Z}}
\newcommand{\Rbb}{\mathbb{R}}
\DeclareMathOperator*{\argmin}{arg\,min}
\begin{document}

\section*{Supplementary Note 7: Independent-batch reproducibility bounds what any perturbation readout can predict}
\label{sup:crossbatch}

\subsection*{Motivation}

Supplementary Figure~S12 reports that the GRQ virtual-perturbation direction does not predict the
real differential-expression response ($\rho\approx0$), while the mean latent-shift direction does
($\rho\approx0.22$). Those comparisons ask how good the \emph{predictor} is. This note asks the
complementary question, which bounds every predictor at once: how reproducible is the
\emph{response itself} in independent experimental batches? A readout cannot be validated against
a target that does not repeat.

\subsection*{Data and design}

We use the published Perturb-seq resource GSE281048 (Jiang, L. et al., \textit{Nat.\ Cell Biol.}
\textbf{27}, 505--517, 2025; Mixscale pipeline), which performs more than
1{,}500 perturbations split across six cell lines and five signalling contexts. Its
\texttt{meta.data} carries an experimental-batch field, \texttt{Batch\_info}, taking the values
\texttt{Rep1} and \texttt{Rep2} (129{,}917 and 106{,}689 cells in the TGFB context; 236{,}606 cells
in total).

\paragraph*{What the two batches are, and what they are not.}
Two features of the design fix the interpretation, and both are reported because they cut in
opposite directions.
\begin{itemize}\itemsep2pt
\item The batches are \emph{not} disjoint cell-line panels: all sixteen \texttt{sample\_ID} labels
appear in both batches (roughly 80\% of each label's cells in \texttt{Rep1}, 65--70\% in
\texttt{Rep2}), so both batches sequence the same material rather than independent biological
preparations. The paper's own description is that the perturbations are ``split across'' the cell
lines. We therefore treat \texttt{Rep1}/\texttt{Rep2} as \emph{technical} batches (experimental
rounds) rather than biological replicates, and read the reproducibility below as a technical
reproducibility.
\item One cell line, BXPC3, is absent from \texttt{Rep1}. This would be a cell-line confound if the
comparison used BXPC3; it does not, because BXPC3 contributes no \texttt{(Rep1, Rep2)} pair. Removing
BXPC3 leaves the group count and the median unchanged ($n=260$, $r=+0.064$).
\end{itemize}

\subsection*{Measurement}

For every (cell line $\times$ target) group with at least twenty cells in \emph{both} batches, the
perturbed cells of each batch are compared with \emph{that batch's own} non-targeting (NT) controls.
Counts are normalised per group to $\log(1+\mathrm{CP10k})$ and the shift is the difference of group
means, so the batch effect is not confounded with the perturbation effect. Reproducibility is the
Pearson correlation between the two batches' shift vectors across all genes.

\subsection*{Controls}

Two controls bracket the measurement, and a third supplies the null.
\begin{itemize}\itemsep2pt
\item \textbf{Positive control (NT vs NT).} The two batches' NT controls agree at $r=+0.9962$ on
average across the five fully shared cell lines (range $+0.9943$ to $+0.9972$). The batch effect on
gene-level pseudobulk is therefore negligible, and the low values below are not a measurement
artefact.
\item \textbf{Mismatched-target null.} Correlating a target's \texttt{Rep1} shift against a
\emph{different} target's \texttt{Rep2} shift, matched on cell line ($2{,}000$ permutations), gives
a median of $+0.0085$ (TGFB) and $+0.0039$ (IFNG).
\item \textbf{Within-batch cross-target floor.} Correlating two different targets' shifts
\emph{inside one batch} gives a median of $+0.2466$ (TGFB).
\end{itemize}

\subsection*{Results}

\begin{table}[htbp]\centering\small
\begin{tabular}{@{}lrrr@{}}
\toprule
& \textbf{TGFB} & \textbf{IFNG} & \textbf{Reference level} \\
\midrule
Same target, across batches            & $+0.0643$ & $+0.0683$ & --- \\
\quad (95\% of values)                 & $n=260$   & $n=345$   & --- \\
NT vs NT, across batches (ceiling)     & $+0.9962$ & $+0.9962$ & measurement ceiling \\
Different target, across batches (null)& $+0.0085$ & $+0.0039$ & $p=6\times10^{-68}$, $6\times10^{-75}$ \\
Different target, within a batch       & $+0.2466$ & ---       & noise floor \\
\bottomrule
\end{tabular}
\caption{\textbf{Cross-batch reproducibility of single-gene perturbation signatures.} The same
target's shift agrees across independent experimental batches at only $r\approx+0.065$, against a
non-targeting ceiling of $+0.996$ and a mismatched-target floor of $+0.009$. The agreement is
significantly above the mismatched null (Mann--Whitney $p=6.2\times10^{-68}$ for TGFB and
$5.7\times10^{-75}$ for IFNG, $n=2{,}000$ permutations), so the signal is real, but it is small:
about $0.4\%$ of the variance at the ceiling is shared across batches. In both contexts fewer than
$1\%$ of groups exceed $r=0.9$.}
\end{table}

Two features of these numbers matter more than their absolute size.

\textbf{First, the same-target agreement is below the within-batch cross-target floor.} Two
\emph{unrelated} perturbations measured in the same batch agree at $+0.247$, while the \emph{same}
perturbation measured in two batches agrees at $+0.064$. Same-batch pairs evidently share a large
common component --- the batch's own state, plus technical structure --- that does not survive
across batches. The batch-stable part of a single-gene signature is therefore smaller than the
part that two arbitrary perturbations happen to share.

\textbf{Second, pathway membership does not raise reproducibility.} Because each context's target
list is a single signalling family, the natural objection is that the reproducible unit is the
\emph{pathway} rather than the gene. We tested this directly, grouping each context's targets into
hand-built functional sets (listed per context in the analysis code; 39/52 TGFB targets and 54/59
IFNG targets covered) and comparing cross-batch correlations within a group against across groups.
Within-group and across-group agreement are indistinguishable: TGFB $+0.0085$ vs $+0.0092$
(difference $-0.0007$, Mann--Whitney $p=0.93$); IFNG $+0.0063$ vs $+0.0058$ (difference $+0.0005$).
The same test inside a single batch also shows no advantage (TGFB \texttt{Rep1}: $+0.2669$ vs
$+0.2605$; \texttt{Rep2}: $+0.2560$ vs $+0.2459$). A small-sample coverage artefact is excluded
because the IFNG grouping was rebuilt after the TGFB grouping was found to cover only 3 of its 59
targets.

\subsection*{Relation to the internal split-half estimate}

An internal estimate obtained by splitting one cell population in half reports a much higher
number: the same groups show a median split-half correlation of about $0.9$ in this resource. The
discrepancy is not a contradiction of the measurement above; it identifies what split-half measures.
Two halves of one population share the library, the batch, and the technical noise of that batch,
so a strong effect drives both halves to the same noisy value and the correlation saturates. This
is also why split-half reproducibility saturates at $1$ in datasets where the underlying
cross-batch agreement is near zero. We therefore report split-half values elsewhere in this work as
an upper bound and a same-batch consistency check, not as evidence that an effect reproduces.

\subsection*{Consequence for perturbation readouts}

The reproducible component of a single-gene perturbation signature is a small fraction of its
variance, and it is not organised by pathway. Any readout --- GRQ, latent-shift, decoder-direct, or
a trained response predictor --- is bounded by this: the ceiling on agreement with a real response
is set by how much of that response is stable across batches, and no method can exceed it. This is
the mechanistic counterpart of the $\rho\approx0$ result for the GRQ direction in Supplementary
Figure~S12, and it is why that null should not be read as a property of the GRQ construction alone.
It also means that target-gene specificity and response prediction can dissociate without either
being defective: a direction can be highly specific to the intended targets and still fail to
predict a response whose reproducible part is small.

\subsection*{Limits}

(i) The batch field is interpreted as technical from design evidence (all \texttt{sample\_ID}
labels present in both batches; the source paper's ``split across'' wording) rather than from an
explicit author definition of the field, which could not be retrieved. If the two batches are in
fact biological replicates, the measured value is a biological reproducibility and is not lower.
(ii) Two contexts were analysed (TGFB, IFNG); the remaining three GSE281048 contexts were not
available locally.
(iii) The pathway groupings are hand-built functional sets by gene name, not database annotations;
they are listed in full in the analysis code so the grouping can be audited, and both contexts
carry their own grouping.
(iv) Correlations are gene-level Pearson on log-normalised pseudobulk without expression
filtering, so low-expression genes contribute noise.
(v) Reproducibility is estimated from two batches only, so no across-batch variance component can
be separated from within-batch variance.

\vspace{10pt}
\noindent\rule{\textwidth}{0.4pt}
\footnotesize\textbf{Source identity.} This note was compiled from \texttt{supplementary/supp\_note7\_crossbatch.tex} in the paper1 repository. SHA256 of that source file with this paragraph removed: \texttt{4c3761851c5d446223f54dc6f4426b70ea366bfa541af32219f03ca8c7768e9c}. Readers can recompute the digest to confirm which revision produced this PDF; the source is distributed with the code release.

\end{document}
